\documentclass[10pt,a4paper]{amsart}
\usepackage[margin=19mm]{geometry}
\usepackage{amsmath,amssymb,mathtools}
\usepackage[scr=rsfs,cal=euler]{mathalfa}
\usepackage{xfrac}
\usepackage[unicode,hidelinks,bookmarks=false]{hyperref}
\newcommand{\ie}{i.e.\ }
\newcommand{\cf}{cf.\ }
\newcommand{\resp}{respectively\ }
\newcommand{\Subsection}{\subsection}
\providecommand{\nobreakdash}{\nobreak\hbox{-}}
\setlength{\emergencystretch}{2em}
\multlinegap0pt
\usepackage{siunitx}
\usepackage{siunitx}
\usepackage{siunitx}
\usepackage{siunitx}
\usepackage{bm}
\usepackage{amssymb}
\usepackage{mathtools}
\usepackage{siunitx}
\usepackage{algorithm}\usepackage{algpseudocode}
\usepackage{xcolor}
\newtheorem{proposition}{Proposition}
\newcommand{\Imiss}{\bm{\mathcal{I}}_{\mathrm{miss}}}
\newcommand{\Io}{\bm{\mathcal{I}}_{\mathrm{obs}}}
\title{Teacher Forcing as Generalized Bayes: Optimization Geometry Mismatch in Switching Surrogates for Chaotic Dynamics}
\author{Andre Herz, Daniel Durstewitz, Georgia Koppe}
\date{Source: arXiv:2604.25904v1 (CC BY 4.0)}
\begin{document}
\maketitle

\noindent\emph{Source and licence.} Teacher Forcing as Generalized Bayes: Optimization Geometry Mismatch in Switching Surrogates for Chaotic Dynamics. Version arXiv:2604.25904v1 (CC BY 4.0), distributed by arXiv under the Creative Commons Attribution 4.0 International licence, http://creativecommons.org/licenses/by/4.0/, as recorded in the arXiv licence metadata for this exact version and shown on the abstract page https://arxiv.org/abs/2604.25904v1. Full licence notice: \texttt{licenses/arxiv-2604.25904-CC-BY-4.0.txt}.\par\smallskip

\noindent\textit{Editorial scope.} Counted unit: the article's proposition prop:s:louis\_missing\_state\_dependent\_gate (source0.tex lines 433--476). The article's complete original proof of it is delivered with it (source0.tex lines 486--512, one \texttt{proof} environment of 27 lines, which the article places away from the statement). No mathematics is written, repaired or re-derived: the statement and the proof are the article's own lines, in the article's own notation, and no retained line is substituted or re-flowed.  No further source-local context is needed: every cross-reference of every retained line resolves inside the two delivered spans.  Nothing was omitted from any retained span, so the package carries no omission record.  No retained line contains a citation, so the wrapper carries no bibliography and no bibliography entry.  No retained line contains a figure environment, an image-inclusion command or drawing code, so no image is delivered and the package declares no asset dependency.  Editorial formatting only, in addition: an AMS \texttt{amsart} preamble that restates the article's own declarations and macros used by the retained lines, namely the environments \texttt{proposition} on separate counters, together with the standard packages those lines need.  The retained environments print in this document as Proposition 1 in order of appearance, whereas the article prints the same items with the numbering of its own full document.  The complete native source of the article as distributed in the arXiv e-print for this exact version is bundled unchanged as \texttt{sources/199341/source0.tex}.
\begin{proposition}[Louis missing-information term under state-dependent gating]
\label{prop:s:louis_missing_state_dependent_gate}
Let $C$ be a latent gate variable and let $\bm{x}$ denote the observed quantity.
For brevity, we write
\begin{equation}
p_{\bm\theta}(\bm x,c):=p_{\bm\theta}(\bm x, C{=}c),\qquad
p_{\bm\theta}(\bm x\mid c):=p_{\bm\theta}(\bm x\mid C{=}c),\qquad
p_{\bm\theta}(c\mid \bm x):=\mathbb{P}_{\bm\theta}(C{=}c\mid \bm x).
\end{equation}
Assume the joint model can be written as
\begin{equation}
p_{\bm\theta}(\bm x,c) \,=\, \pi(c\mid \bm x)\,p_{\bm\theta}(\bm x\mid c),\qquad c\in\{0,1\},
\end{equation}
where $\pi(c\mid \bm x)$ is a (possibly $\bm x$-dependent) gating term that is \emph{independent of $\bm\theta$} for the parameter block of interest.
Define, for each $c\in\{0,1\}$,
\begin{equation}
\bm s_c(\bm\theta) := \nabla_{\bm\theta}\log p_{\bm\theta}(\bm x\mid c),
\qquad
\bm{\mathcal I}_c(\bm\theta) := -\nabla_{\bm\theta}^2\log p_{\bm\theta}(\bm x\mid c).
\end{equation}
Let
\begin{equation}
p := p_{\bm\theta}(1\mid \bm x)
= \frac{\pi(1\mid \bm x)\,p_{\bm\theta}(\bm x\mid 1)}{\pi(0\mid \bm x)\,p_{\bm\theta}(\bm x\mid 0)+\pi(1\mid \bm x)\,p_{\bm\theta}(\bm x\mid 1)}.
\end{equation}
Then the observed information for the marginal likelihood $p_{\bm\theta}(\bm x)$ satisfies
\begin{equation}
\label{eq:s:2switch_state_dependent}
\Io(\bm\theta)
=
p\,\bm{\mathcal{I}}_1(\bm\theta) + (1-p)\,\bm{\mathcal{I}}_0(\bm\theta)
- \Imiss(\bm\theta).
\end{equation}
\begin{equation}
\label{eq:s:2switch_state_dependent_miss}
\Imiss(\bm\theta)
:=
p(1-p)\big(\bm{s}_1(\bm\theta)-\bm{s}_0(\bm\theta)\big)\big(\bm{s}_1(\bm\theta)-\bm{s}_0(\bm\theta)\big)^{\mathsf T}
\succeq 0.
\end{equation}
In particular, the missing-information term $\Imiss(\bm\theta)$ is large when (i) the posterior is ambiguous
($p(1-p)$ is maximized at $p=1/2$) and (ii) the two explanations disagree in score
(large $\|\bm{s}_1(\bm\theta)-\bm{s}_0(\bm\theta)\|$).
\end{proposition}

\begin{proof}[Proof of Proposition~\ref{prop:s:louis_missing_state_dependent_gate}]
Because $\pi(c\mid \bm x)$ is $\bm\theta$-independent, we have
\begin{equation}
\log p_{\bm\theta}(\bm x,c)=\log \pi(c\mid \bm x)+\log p_{\bm\theta}(\bm x\mid c).
\end{equation}
Therefore the complete-data score and complete-data negative Hessian for the parameter block $\bm\theta$ are
\begin{equation}
\nabla_{\bm\theta}\log p_{\bm\theta}(\bm x,c)=\bm s_c(\bm\theta),
\qquad
-\nabla_{\bm\theta}^2\log p_{\bm\theta}(\bm x,c)=\bm{\mathcal I}_c(\bm\theta).
\end{equation}
Louis' identity for this two-point latent gate can be written as
\begin{align}
\Io(\bm\theta)
&=\sum_{c\in\{0,1\}}p_{\bm\theta}(c\mid \bm x)\,\bm{\mathcal I}_c(\bm\theta)\notag\\
&\quad-\Bigg(
\sum_{c\in\{0,1\}}p_{\bm\theta}(c\mid \bm x)\,\bm s_c(\bm\theta)\bm s_c(\bm\theta)^{\mathsf T}
-\Big(\sum_{c\in\{0,1\}}p_{\bm\theta}(c\mid \bm x)\,\bm s_c(\bm\theta)\Big)\Big(\sum_{c\in\{0,1\}}p_{\bm\theta}(c\mid \bm x)\,\bm s_c(\bm\theta)\Big)^{\mathsf T}
\Bigg).
\end{align}
Writing $p:=p_{\bm\theta}(1\mid \bm x)$ and simplifying yields
\begin{equation}
\Io(\bm\theta)=p\,\bm{\mathcal I}_1(\bm\theta)+(1-p)\,\bm{\mathcal I}_0(\bm\theta)
-p(1-p)\big(\bm s_1(\bm\theta)-\bm s_0(\bm\theta)\big)\big(\bm s_1(\bm\theta)-\bm s_0(\bm\theta)\big)^{\mathsf T},
\end{equation}
which is Eq.~\eqref{eq:s:2switch_state_dependent} with the missing-information term shown in Eq.~\eqref{eq:s:2switch_state_dependent_miss}.
\end{proof}

\end{document}
